\title{\bfseries\fontsize{18}{24}\selectfont 带符号循环图的反例与通量相结构}
\ifdefined\TargetAAnonymous
  \author{\fontsize{12}{18}\selectfont 匿名}
\else
  \author{\fontsize{12}{18}\selectfont [AUTHOR NAME]}
\fi
\date{}
\maketitle

\ifdefined\TargetAAnonymous\else
\begin{center}
\small
\mbox{[AFFILIATION]}\\
\mbox{[DEPARTMENT]}, \mbox{[INSTITUTION]}\\
\mbox{[CITY]}, \mbox{[COUNTRY]}\\[3pt]
通讯作者：[CORRESPONDING AUTHOR]\\
电子邮箱：[EMAIL]\quad ORCID：[ORCID]
\end{center}
\fi

\begin{abstract}
对于带符号循环图 \(C_n(1,2)\)，近期一个猜想断言：当 \(n\geq 8\) 为偶数时，
具有负 Hamilton 圈和乐且三角形通量交替的等价类总能使谱半径达到最小。
本文否定该猜想并确定其最小失效阶数。精确穷举计算验证了所有偶数
\(8\leq n\leq 30\) 上的猜想下界，而 \(n=32\) 时一个显式赋号具有严格更小的
谱半径；因此有限阶最小性结论是计算机辅助的。该反例由一个周期八通量模式
生成，并在所有满足 \(8\mid n\) 且 \(n\geq 32\) 的阶数上产生反例。Floquet
约化给出平方能带的精确四次方程以及周期八的锐值
\(\eta=4+\sqrt{10+2\sqrt{5}}\)。随后我们解释这一相位背后的机制：带符号算子
平方后，负通量恰好消去所有相关的奇位移耦合，而正通量产生局域缺陷。闭步
矩给出完整的周期八三分律：对径缺陷对恰为谱平方半径小于八的相位；全负相位
的谱平方半径等于八；其余每个相位均大于八。对于任意周期，我们导出前三个
精确矩，并得到缺陷密度和局部聚集的必要不等式。最后，一项完整的精确分类
表明，在本原 Hamilton 规范周期至多为十六的周期相位中，除自然算子等价外，
该周期八相位是唯一极小元。
\end{abstract}

\noindent\textbf{关键词：} [KEYWORDS TO CONFIRM]\par
\noindent\textbf{MSC 2020：} [MSC CODES TO CONFIRM]
